\begin{center}
\textbf{\fontsize{16}{32}\selectfont ABSTRACT}
\end{center}
\vspace{1em}

\paragraph*{}Vibration-based condition monitoring serves as a vital cornerstone for ensuring the operational safety, mechanical integrity, and predictive maintenance of rotating electromechanical machinery. This project presents the design, embedded implementation, and empirical validation of \textbf{VibeGuard}, a low-cost, edge-resident vibration condition monitoring system. VibeGuard acquires triaxial acceleration data at an 800\,Hz Output Data Rate (ODR) from an Analog Devices ADXL345 MEMS sensor interfaced with a 32-bit dual-core 7Semi ESP32-DEVKIT-E microcontroller over a deterministic 5\,MHz 4-wire SPI bus. The microcontroller executes an autonomous digital signal processing pipeline directly at the edge, performing continuous in-place zero-mean DC gravity cancellation, 3-axis Vector RMS ($V_{\text{RMS}}$) feature extraction, and temporal $K$-of-$M$ persistence filtering ($K=4$ violations out of $M=5$ consecutive 320\,ms windows) without relying on external cloud dependencies.

\paragraph*{}To validate the system, an experimental testbed comprising a 12\,V DC micro metal gear motor running at 600\,RPM (10.0\,Hz rotational shaft frequency) and an eccentric 3D-printed rotor arm was evaluated across multiple controlled operational stages. Under healthy baseline conditions (Stage 0), the system established an autonomous reference with $V_{\text{RMS}} \approx 0.013$--$0.060\,g$, illuminating a solid Green LED. When a moderate unbalance mass was introduced (Stage 1: bolt only, $m = 1.2\,\text{g}$), the composite vibration rose to $0.82$--$0.86\,g$, decisively asserting a solid Yellow warning state with zero spurious alarm transitions. Under severe mechanical unbalance (Stage 2/3: bolt plus locknuts, $m = 1.85\,\text{g}$), dynamic vibration surged to $1.079\,g$ (a 423\% increase), triggering a high-visibility 6.25\,Hz Fast Red Strobe alert within 1.28\,seconds. On-demand 256-point Hann-windowed Fast Fourier Transform (FFT) analysis confirmed a 16-fold spectral energy concentration centered precisely at the 10.0\,Hz synchronous 1X mechanical unbalance peak. Screwdriver impulse shock testing confirmed complete transient rejection ($V = 1 < 4$, zero false alarms). Telemetry was streamed to a local web-based dashboard for real-time visualization. Consuming under 1.5\,W, VibeGuard demonstrates an explainable, robust, and cost-effective paradigm for industrial asset protection and embedded engineering education.
\clearpage
